% !TeX root = ../../Book5.tex
% 纲要标题：### 20.4 BGG 方法
% 纲要位置：工作记录/计划纲要.md，第 4381 行。

\section{BGG 方法}
\label{b5:sec:2004}

% 纲要范围：
%
% * BGG 消解的典型版本
% * Ext 与最高权范畴
% * 几何表示论所需的进一步工具
%
% ---
%

取最高权模的简单商时，要商去较低最高权生成的子模。若出现多个这样的子模，还须考虑它们的交及相互关系。BGG 消解通过一个正合复形逐次描述这些关系。

\subsection{BGG 消解的典型版本}
\label{b5:subsec:2004:01}

\begin{theorem}[BGG 消解]\label{b5:thm:bgg-resolution}
设 \(\mathfrak g\) 为有限维复半单 Lie 代数，固定 Cartan 子代数 \(\mathfrak h\) 及正根集合 \(\Phi^+\)，\(\lambda\in\mathfrak h^*\) 为相应的支配整权，\(W\) 为 Weyl 群，\(\rho=\dfrac12\sum_{\alpha\in\Phi^+}\alpha\)。记 \(M(\mu)\) 为最高权 \(\mu\in\mathfrak h^*\) 的 Verma 模，\(L(\mu)\) 为其唯一简单商，\(\ell(w)\) 为 \(w\in W\) 的简单反射长度，\(w_0\) 为最长元，\(N=\ell(w_0)\)。对 \(0\le j\le N\)，记
\[
D_j=\bigoplus_{\ell(w)=j}M(w\mathbin{\cdot}\lambda),
\qquad w\mathbin{\cdot}\lambda=w(\lambda+\rho)-\rho.
\]
可以选择模同态 \(D_j\to D_{j-1}\)，使
\[
0\longrightarrow D_N\longrightarrow\cdots\longrightarrow D_1
\longrightarrow D_0\longrightarrow L(\lambda)\longrightarrow0
\]
正合。它称为\term[BGG-xiaojie]{BGG 消解}[Bernstein--Gelfand--Gelfand resolution]。
\end{theorem}
\symidx{\(w\mathbin{\cdot}\lambda\)：Weyl 群在权上的移位作用}

\begin{proofsketch}
先取 \(\lambda=0\)。令 \(\mathfrak b=\mathfrak h\oplus\mathfrak n_+\)，以 \(\mathfrak b\) 的伴随作用把 \(\mathfrak g/\mathfrak b\) 看成 \(\mathfrak b\) 模。考虑相对 Chevalley--Eilenberg 复形的各项
\[
 E_j=U(\mathfrak g)\otimes_{U(\mathfrak b)}
 \bigwedge^j(\mathfrak g/\mathfrak b),
 \qquad 0\le j\le N.
\]
以 \(\bar x=x+\mathfrak b\) 记商空间中的类。前两次微分具体为
\[
\begin{aligned}
d_1(u\otimes\bar x)&=ux\otimes1,\\
d_2(u\otimes\bar x\wedge\bar y)
&=ux\otimes\bar y-uy\otimes\bar x-u\otimes\overline{[x,y]}.
\end{aligned}
\]
括号取在 \(\mathfrak g\) 中，再取商类；\(\mathfrak g/\mathfrak b\) 本身并非商 Lie 代数。若把 \(x\) 换成 \(x+b\)、\(b\in\mathfrak b\)，第一式的差为 \(ub\otimes1=u\otimes b\cdot1=0\)，第二式的差为
\[
ub\otimes\bar y-u\otimes\overline{[b,y]}=0,
\]
因为张量积是在 \(U(\mathfrak b)\) 上取的，且 \(b\cdot\bar y=\overline{[b,y]}\)。因此提升改变造成的作用项与括号项恰好抵消。一般微分同样由删去一个外幂因子的作用项与两两括号项组成，Jacobi 恒等式使相邻微分复合为零。PBW 分解 \(\mathfrak g=\mathfrak n_-\oplus\mathfrak b\) 将此复形作为向量空间识别为 \(\mathfrak n_-\) 的标准消解；给 \(U(\mathfrak n_-)\) 取 PBW 滤过后，其分次复形是对称代数的 Koszul 消解，所以除了末端的平凡模外没有同调。

这里还要使用广义中心特征标的分解。对 \(M\in\mathcal O\)，以 \(M_\chi\) 记具有广义中心特征标 \(\chi\) 的向量组成的子模，则
\[
 M=\bigoplus_\chi M_\chi,
 \qquad \operatorname{pr}_\chi:M\longmapsto M_\chi
\]
只有有限多个非零项，且 \(\operatorname{pr}_\chi\) 为正合函子。其依据是中心保持每个有限维权空间，交换算子的广义共同特征空间分解与模同态相容；有限生成性又将出现的特征标限制为有限多个。对短正合列逐权作这一分解，便得到投影的正合性。这一分解见\cite[Section 1.12]{Humphreys2008BGGCategoryO}。\footnote{本节关于中心分解、张量滤过与 BGG 消解证明的定位按 Humphreys 的 2008 年作者稿核对，所引节号与 AMS 版目录中的同名各节一致。}

这些 \(E_j\) 具有 Verma 滤过，投影 \(\operatorname{pr}_{\chi_0}\) 因而仍给出正合复形。\(\bigwedge^j(\mathfrak g/\mathfrak b)\) 的权是 \(j\) 个不同负根的和；外幂权引理说明，其中与零权处于同一移位 \(W\) 轨道的，恰为 \(w\mathbin{\cdot}0=w\rho-\rho\)、\(\ell(w)=j\)，各出现一次。于是投影后的第 \(j\) 项具有每个 \(M(w\mathbin{\cdot}0)\) 恰一次的 Verma 滤过。为使它成为直和，需要 Verma 扩张判据的下述推论：若 \(\mu\) 为支配整权，\(x,y\in W\) 满足 \(\ell(x)=\ell(y)\)，则
\[
 \operatorname{Ext}^1_{\mathcal O}
 \bigl(M(x\mathbin{\cdot}\mu),M(y\mathbin{\cdot}\mu)\bigr)=0.
\]
事实上，非零扩张要求 \(y<x\) 成立于 Bruhat 次序，因而 \(\ell(y)<\ell(x)\)；该判据见\cite[Section 6.5]{Humphreys2008BGGCategoryO}。逐次分裂滤过中的短正合列，就得到 \(\lambda=0\) 时的 \(D_j\)。

例如取 \(\mathfrak{sl}_3\)，简单根记为 \(\alpha,\beta\)。商空间 \(\mathfrak g/\mathfrak b\) 的三个权为 \(-\alpha,-\beta,-\alpha-\beta\)。外幂次数零至三的权依次是
\[
\begin{aligned}
j=0:&\quad 0,\\
j=1:&\quad-\alpha,\ -\beta,\ -\alpha-\beta,\\
j=2:&\quad-\alpha-\beta,\ -2\alpha-\beta,\ -\alpha-2\beta,\\
j=3:&\quad-2\alpha-2\beta.
\end{aligned}
\]
其中零权的移位轨道为
\(0,-\alpha,-\beta,-2\alpha-\beta,-\alpha-2\beta,-2\alpha-2\beta\)。所以中心投影在次数一删去 \(-\alpha-\beta\) 对应的 Verma 因子，在次数二也删去同权的因子，保留的因子数依次为 \(1,2,2,1\)。这首先是投影后的 Verma 滤过；还要调用刚才的同长度 \(\operatorname{Ext}^1\) 消失，才可把每一项写成这些 Verma 模的直和。中心投影本身不提供这一步分裂。

一般支配整权 \(\lambda\) 可从这一复形出发，先逐项张量有限维模 \(L(\lambda)\)，再施行 \(\operatorname{pr}_{\chi_\lambda}\)。这两步均正合，末端变为 \(L(\lambda)\)。张量与诱导的标准恒等式给出 \(M(w\mathbin{\cdot}0)\otimes L(\lambda)\) 的 Verma 滤过：因子为
\[
 M(w\rho-\rho+\nu),
 \qquad \text{重数为}\;\dim L(\lambda)_\nu,
\]
其中 \(\nu\) 遍历 \(L(\lambda)\) 的权；见\cite[Section 3.6]{Humphreys2008BGGCategoryO}。现在核对投影保留哪些因子。由\cref{b5:thm:harish-chandra}，若这个因子的中心特征标是 \(\chi_\lambda\)，就存在 \(t\in W\) 使
\[
 t(w\rho+\nu)=\rho+\lambda.
\]
记 \(Q_+\) 为简单根的非负整数组合。因 \(\rho\) 支配，\(\rho-tw\rho\in Q_+\)；又因有限维模的权在 \(W\) 下不变且不高于最高权，\(\lambda-t\nu\in Q_+\)。上式使这两项之和为零，所以两项均为零。\(\rho\) 的稳定子为平凡群，故 \(t=w^{-1}\)，从而 \(\nu=w\lambda\)。这个极端权的重数为一。因此投影只留下 \(M(w\mathbin{\cdot}\lambda)\) 一个 Verma 因子。投影与有限直和相容，所得第 \(j\) 项正是 \(D_j\)，微分由原复形诱导并保持正合。

相对微分的良定义及滤过消解核验、外幂权引理和 Verma 扩张判据的证明见\cite[Sections 6.2--6.5]{Humphreys2008BGGCategoryO}。\cite[Section 8.4, Theorem 8.30]{Kirillov2008LieGroups}给出了同一定理的陈述，但未展开证明。
\end{proofsketch}

下面在秩一情形给出具体映射，并直接证明正合性。

\begin{proposition}\label{b5:prop:bgg-sl2-explicit}
设 \(\mathfrak g=\mathfrak{sl}_2(\mathbb C)\)，\(m\in\mathbb Z_{\ge0}\)，记 \(M(\mu)\) 为 \(\mathfrak g\) 的最高权 \(\mu\) 的 Verma 模，\(L_m\) 为最高权 \(m\)、维数 \(m+1\) 的不可约 \(\mathfrak g\) 模。令 \(v,u\) 分别为 \(M(m),M(-m-2)\) 的非零最高权生成元，\(f=E_{21}\)，其中 \(E_{21}\) 为矩阵单位。则下式给出一个不分裂的短正合列，\(\iota\) 按所有整数 \(j\ge0\) 时的公式定义：
\[
0\longrightarrow M(-m-2)\overset{\iota}{\longrightarrow}
M(m)\longrightarrow L_m\longrightarrow0,
\qquad \iota(f^ju)=f^{j+m+1}v,
\]
\end{proposition}
\begin{proof}
向量 \(f^{m+1}v\) 的权为 \(-m-2\)，且
\(ef^{m+1}v=(m+1)(m-m)f^mv=0\)。Verma 泛性质给出 \(\iota\)；PBW 基表明所有 \(f^{j+m+1}v\) 线性无关，所以 \(\iota\) 单射。商模以 \(v,fv,\ldots,f^mv\) 为基，\(e,f,h\) 的作用正是有限维简单模 \(L_m\) 的标准公式，故商为 \(L_m\)。

若正合列分裂，商中最高向量的像必须是 \(M(m)\) 中权 \(m\) 的向量，因此是 \(v\) 的非零倍数。这个向量生成整个 \(M(m)\)，却在 \(L_m\) 中应被 \(f^{m+1}\) 零化，矛盾。因此它不分裂。
\end{proof}

在一般消解中，逐权取维数仍合法，因为每个权空间有限维。欧拉交错和给出
\[
\operatorname{ch}L(\lambda)
=\sum_{w\in W}(-1)^{\ell(w)}\operatorname{ch}M(w\mathbin{\cdot}\lambda).
\]
再用 PBW 的
\(\operatorname{ch}M(\mu)=e^\mu\prod_{\alpha>0}(1-e^{-\alpha})^{-1}\)，即可恢复 Weyl 公式。这个推导从 BGG 消解解释前章的公式。

\subsection{Ext 与最高权范畴}
\label{b5:subsec:2004:02}

\begin{definition}\label{b5:def:yoneda-ext-o}
设 \(\mathfrak g\) 为有限维复半单 Lie 代数，固定 Cartan 子代数及正根集合，令 \(\mathcal O\) 为相应的范畴 O，\(M,N\in\mathcal O\)。中间项 \(E\) 也在 \(\mathcal O\) 内的短正合列
\(0\to N\to E\to M\to0\)，按在两端为恒等的交换图取等价类，并用拉回、推出定义 Baer 加法，所得群记作 \(\operatorname{Ext}^1_{\mathcal O}(M,N)\)。零类恰为分裂列。
\end{definition}
\symidx{\(\operatorname{Ext}^1_{\mathcal O}(M,N)\)：范畴 O 中短正合扩张的等价类群}

\cref{b5:prop:bgg-sl2-explicit}因此给出
\(\operatorname{Ext}^1_{\mathcal O}\bigl(L_m,M(-m-2)\bigr)\ne0\)。这里不能把 BGG 消解误称为投射消解：它的各项是 Verma 模，而 Verma 模并非都投射。

\begin{example}\label{b5:exm:o-costandard-sl2}
在 \(M(m)\) 的每个一维权空间上取对偶，令 \(\varphi_j\) 与 \(f^jv\) 对偶，取其代数直和 \(D M(m)=\bigoplus_{j\ge0}\mathbb C\varphi_j\)。具体地，令 \(\tau(e)=f\)、\(\tau(f)=e\)、\(\tau(h)=h\)，并按 \(\tau(uv)=\tau(v)\tau(u)\) 延拓为包络代数的反自同构，定义
\[
(x\varphi)(v)=\varphi(\tau(x)v).
\]
这个作用没有普通对偶中的负号：它是借 \(\tau\) 反转乘法，故 \(h\) 权保持为原来的权。直和中的泛函只在有限多个权空间上非零；整个代数对偶 \(M(m)^*\cong\prod_{j\ge0}\mathbb C\varphi_j\) 则允许在无限多个权空间上取非零值。因此这里取的是受限对偶，而非整个代数对偶。按定义逐基计算得到
\[
h\varphi_j=(m-2j)\varphi_j,\qquad
e\varphi_j=\varphi_{j-1},\qquad
f\varphi_j=(j+1)(m-j)\varphi_{j+1},
\]
其中 \(\varphi_{-1}=0\)。这些公式也可以直接验证三个 \(\mathfrak{sl}_2\) 括号。
若 \(m\ge0\)，前 \(m+1\) 个向量生成子模 \(L_m\)；商的最高权为 \(-m-2\)，后续 \(f\) 作用均非零，故经逐项缩放基后商同构于 \(M(-m-2)\)。中间模由 \(\varphi_{m+1}\) 生成，具有有限维权空间且 \(e\) 局部幂零，所以属于 \(\mathcal O\)。得到
\[
0\longrightarrow L_m\longrightarrow D M(m)
\longrightarrow M(-m-2)\longrightarrow0.
\]
该列不分裂：商的最高向量只有 \(\varphi_{m+1}\) 的标量倍提升能保持权，而
\(e\varphi_{m+1}=\varphi_m\ne0\)。于是 \(M(-m-2)\) 不是 \(\mathcal O\) 中的投射对象。
\end{example}

最高权范畴把简单模、标准模和投射模放在同一个偏序下。一个常用的有限版本如下：设有限长度的 \(\mathbb C\) 线性交换范畴具有有限维同态空间，有有限多个简单对象 \(L(\lambda)\)，由有限偏序集标记，并有投射覆盖 \(P(\lambda)\)。若可选择标准对象 \(\Delta(\lambda)\)，使其顶为 \(L(\lambda)\)、其余合成因子参数严格小于 \(\lambda\)，且
\(\operatorname{End}\Delta(\lambda)=\mathbb C\)，同时
\(P(\lambda)\twoheadrightarrow\Delta(\lambda)\)
的核有参数严格大于 \(\lambda\) 的标准对象滤过，就得到一种\term[zuigaoquan-fanchou]{最高权范畴}[highest weight category]结构。这个定义还要求投射覆盖与上述滤过确实存在；简单模有最高权本身并不足以推出这些条件。

\ProseParagraphBreak
可以直接看见投射性如何依赖范畴范围。对 \(\mathfrak{sl}_2\)，在权只允许为
\(m,m-2,m-4,\ldots\) 的 \(\mathcal O\) 全子范畴中，
\[
\operatorname{Hom}\bigl(M(m),N\bigr)\cong N_m.
\]
因为 \(N\) 中没有权 \(m+2\)，任何权 \(m\) 的向量都被 \(e\) 零化，Verma 泛性质便给出这个同构。取一个固定权空间是正合函子，故 \(M(m)\) 在这个受限范畴中投射。这个计算与上面的非投射例并不矛盾：投射性的判断始终包含允许哪些对象这一条件。

\ProseParagraphBreak
一般 \(\mathcal O\) 的分块与最高权结构还需要有限长度、投射覆盖和标准滤过理论，本节暂不展开。计算高阶 \(\operatorname{Ext}\) 时，应使用相应范畴中的投射或内射消解，或先证明替代消解的各项相对于目标函子为零调对象。BGG 消解的正合性本身尚不足以说明对它施加 \(\operatorname{Hom}\) 就能计算 \(\operatorname{Ext}\)。

\subsection{几何表示论所需的进一步工具}
\label{b5:subsec:2004:03}

几何方法从旗簇上的局部对象构造表示。在 \(\mathfrak{sl}_2\) 情形，旗簇就是 \(\mathbb P^1\)。次数为 \(m\ge0\) 的齐次多项式是线丛 \(\mathcal O(m)\) 的整体截面，而算子
\[
e=X\frac{\partial}{\partial Y},\qquad
f=Y\frac{\partial}{\partial X},\qquad
h=X\frac{\partial}{\partial X}-Y\frac{\partial}{\partial Y}
\]
使它成为 \(L_m\)。在坐标 \(z=Y/X\) 下，这些算子变成
\(e=\partial_z\)、\(h=m-2z\partial_z\)、\(f=mz-z^2\partial_z\)。
\appref{b5:app:C}将从坐标变换核验这些局部微分算子以及相应中心关系。

\ProseParagraphBreak
高维旗簇的 Schubert 分解\footnote{舒伯特（Hermann Cäsar Hannibal Schubert，1848-1911），德国数学家，开创了代数几何中的枚举几何研究。}按 Weyl 群编号，但胞腔闭包可能奇异。一般简单模的重数问题因而需要层的导出函子、交叉上同调和微分算子模等工具；光滑胞腔与其可能奇异的闭包也须区分。关于 BGG 与旗簇、几何表示的关系，可参见\cite[Section 8.4 and Further reading, Geometric representation theory]{Kirillov2008LieGroups}。这些几何方法进一步研究表示的构造与重数。
